%==============================================================================
% figures/tikz/usecase-actors.tex
% Human actor taxonomy for Bayan — redesigned v4 (2026-10) per UML standards.
%
% CHANGES vs v3:
%   - AI agents REMOVED entirely: per UML, the five agents are internal system
%     components (nodes inside the boundary), NOT actors. They are modeled in
%     Section 3.5 (Agent Modeling) instead.
%   - Two-level ORGANISATIONAL taxonomy of the four HUMAN actors:
%       Level 1: Platform Administrator (PA)
%       Level 2: Tenant Users (TU) — abstract actor generalising 3 concrete
%                roles: Tenant IT Administrator, Enterprise Manager,
%                Enterprise Employee (UML generalisation, hollow triangle).
%   - Numbered badges (1-4) match the rows of tab:use-case-actors.
% Every Arabic label is wrapped in \textarabic{} inside bidi's LTR
% environment (same proven fix as the Gantt charts) so word order renders
% correctly. Card order runs right-to-left to match Arabic reading.
%==============================================================================
\newsavebox{\ActorsBox}%
\begin{lrbox}{\ActorsBox}%
\begin{LTR}%
\begin{tikzpicture}[
  font=\small,
  stickman/.pic={
    \draw[line width=1.1pt, #1] (0,0.4) circle (0.2);
    \draw[line width=1.1pt, #1] (0,0.2) -- (0,-0.3);
    \draw[line width=1.1pt, #1] (-0.3,0.06) -- (0,-0.06) -- (0.3,0.06);
    \draw[line width=1.1pt, #1] (0,-0.3) -- (-0.22,-0.75);
    \draw[line width=1.1pt, #1] (0,-0.3) -- (0.22,-0.75);
  },
  server/.pic={
    % Stacked server racks icon (Platform Administrator)
    \draw[line width=1.05pt, #1, rounded corners=2pt] (-0.34,0.14) rectangle (0.34,0.44);
    \draw[line width=1.05pt, #1, rounded corners=2pt] (-0.34,-0.20) rectangle (0.34,0.10);
    \draw[line width=1.05pt, #1, rounded corners=2pt] (-0.34,-0.54) rectangle (0.34,-0.24);
    \fill[#1] (-0.20,0.29) circle (0.035);
    \fill[#1] (-0.20,-0.05) circle (0.035);
    \fill[#1] (-0.20,-0.39) circle (0.035);
    \draw[line width=0.9pt, #1] (0.04,0.29) -- (0.24,0.29);
    \draw[line width=0.9pt, #1] (0.04,-0.05) -- (0.24,-0.05);
    \draw[line width=0.9pt, #1] (0.04,-0.39) -- (0.24,-0.39);
  },
  actorcard/.style={
    rounded corners=6pt, draw=#1!75, fill=white, line width=0.9pt,
    minimum width=3.7cm, minimum height=3.05cm
  },
  widecard/.style={
    rounded corners=6pt, draw=#1!75, fill=white, line width=0.9pt,
    minimum width=10.6cm, minimum height=1.62cm
  },
  actornum/.style={
    circle, fill=#1, text=white, font=\bfseries\scriptsize,
    inner sep=0pt, minimum size=0.42cm
  },
  actorname/.style={
    align=center, font=\bfseries\scriptsize, text=kaudark,
    inner sep=0pt
  },
  actoreng/.style={
    align=center, font=\tiny, text=kaudark!70, inner sep=0pt
  },
  actordesc/.style={
    align=center, font=\tiny, text=kaudark!85, inner sep=0pt,
    text width=3.25cm
  },
  bandpill/.style={
    rounded corners=4pt, fill=#1, text=white, inner sep=5pt,
    font=\bfseries\scriptsize
  },
  gen/.style={draw=kaudark!65, line width=0.85pt}
]

  % ============ LEVEL 1: Platform Administrator (PA) ============
  \fill[kauprimary!6, rounded corners=10pt] (0.15,0.62) rectangle (15.85,-2.30);
  \node[bandpill=kauprimary] at (8.0,0.55)
    {\textarabic{المستوى الأول: مسؤول المنصة \EN{Platform Administrator (PA)}}};

  \draw[widecard=kauprimary] (2.7,-0.85) rectangle (13.3,-2.0);
  \pic at (4.05,-1.42) {server={kaudark}};
  \node[actornum=kauprimary] at (12.9,-1.08) {\EN{1}};
  \node[actorname, anchor=west] at (5.35,-1.18)
    {\textarabic{مسؤول المنصة \EN{(Platform Administrator — PA)}}};
  \node[actordesc, anchor=west, text width=7.4cm] at (5.35,-1.72)
    {\textarabic{تشغيل المنصة تقنياً: التهيئة الذكية للمؤسسات، وإدارة محرك الذكاء الاصطناعي ومزودي النماذج (\EN{Ollama} / السحابة)، ومراقبة صحة الحاويات والنظام}};

  % ============ LEVEL 2: Tenant Users (TU) — abstract + 3 roles ============
  \fill[kausecondary!6, rounded corners=10pt] (0.15,-2.70) rectangle (15.85,-9.42);
  \node[bandpill=kausecondary] at (8.0,-2.77)
    {\textarabic{المستوى الثاني: مستخدمو المؤسسة \EN{Tenant Users (TU)} --- فاعل مجرد يتفرع بالتعميم إلى ثلاثة أدوار}};

  % Abstract parent chip (dashed = abstract actor)
  \node[rounded corners=5pt, draw=kausecondary, dashed, line width=0.9pt,
        fill=white, inner sep=5pt, font=\scriptsize\bfseries, text=kaudark]
    (tuparent) at (8.0,-3.62)
    {\textarabic{مستخدم المؤسسة \EN{(Tenant Users — TU)}}\;\textit{\tiny\EN{\{abstract\}}}};

  % Generalisation bus: 3 children -> TU parent (hollow triangle at parent)
  \draw[gen] (13.1,-4.95) -- (13.1,-4.35) -- (8.0,-4.35);
  \draw[gen] (8.0,-4.95) -- (8.0,-4.35);
  \draw[gen] (2.9,-4.95) -- (2.9,-4.35) -- (8.0,-4.35);
  \draw[gen, -{Triangle[open, length=2.4mm, width=2.2mm]}] (8.0,-4.35) -- (8.0,-3.98);
  \node[font=\tiny, text=kaudark!70] at (10.35,-4.55)
    {\textarabic{تعميم \EN{(Generalization)}}};

  % Role 1 (rightmost, first read): Tenant IT Administrator
  \draw[actorcard=kausecondary] (11.25,-4.95) rectangle (14.95,-8.0);
  \pic at (13.1,-5.75) {stickman={kaudark}};
  \node[actornum=kausecondary] at (14.6,-5.25) {\EN{2}};
  \node[actorname] at (13.1,-6.68) {\textarabic{مسؤول تقنية المعلومات}\\[-1pt]\textarabic{بالمؤسسة}};
  \node[actoreng]  at (13.1,-7.24) {\EN{Tenant IT Administrator}};
  \node[actordesc] at (13.1,-7.68) {\textarabic{إدارة حسابات المستخدمين والأدوار (\EN{RBAC})، وربط مصادر البيانات وتصنيفها الأمني، ومراجعة سجل التدقيق}};

  % Role 2: Enterprise Manager
  \draw[actorcard=kausecondary] (6.15,-4.95) rectangle (9.85,-8.0);
  \pic at (8.0,-5.75) {stickman={kaudark}};
  \node[actornum=kausecondary] at (9.5,-5.25) {\EN{3}};
  \node[actorname] at (8.0,-6.68) {\textarabic{مدير المؤسسة}};
  \node[actoreng]  at (8.0,-7.24) {\EN{Enterprise Manager}};
  \node[actordesc] at (8.0,-7.68) {\textarabic{استعلامات البيانات المهيكلة والبحث في اللوائح، وطلب التقارير التحليلية واكتشاف القيم الشاذة}};

  % Role 3: Enterprise Employee
  \draw[actorcard=kausecondary] (1.05,-4.95) rectangle (4.75,-8.0);
  \pic at (2.9,-5.75) {stickman={kaudark}};
  \node[actornum=kausecondary] at (4.4,-5.25) {\EN{4}};
  \node[actorname] at (2.9,-6.68) {\textarabic{موظف المؤسسة}};
  \node[actoreng]  at (2.9,-7.24) {\EN{Enterprise Employee}};
  \node[actordesc] at (2.9,-7.68) {\textarabic{الاستعلامات التشغيلية اليومية ضمن حدود الصلاحية، والبحث الدلالي في اللوائح والسياسات}};

  % Footnote: agents are internal components (Section 3.5)
  \node[rounded corners=4pt, draw=dividerdim!80, dashed, fill=white,
        inner sep=5pt, font=\tiny, text=kaudark!85, align=center] at (8.0,-8.95)
    {\textarabic{وكلاء الذكاء الاصطناعي الخمسة ليست فاعلين؛ فهي مكونات داخلية ضمن حدود النظام تُنمذَج في القسم \EN{3.5} (نمذجة منظومة الوكلاء)}};

\end{tikzpicture}%
\end{LTR}%
\end{lrbox}%
\resizebox{\textwidth}{!}{\usebox{\ActorsBox}}%
